\begin{mistake}{2026-06-30}{01}

\begin{problemcontent}
设函数 \(f(x)\) 在 \([0,3]\) 上连续，在 \((0,3)\) 内存在二阶导数，且
\[
2f(0)=\int_0^2 f(x)\,dx=f(2)+f(3).
\]
证明：
\begin{enumerate}
  \item 存在 \(\eta\in(0,2)\)，使得 \(f(\eta)=f(0)\)；
  \item 存在 \(\xi\in(0,3)\)，使得 \(f'(\xi)=0\)。
\end{enumerate}
\end{problemcontent}

\begin{solutioncontent}
令
\[
g(x)=f(x)-f(0).
\]
由 \(f(x)\) 在 \([0,3]\) 上连续，知 \(g(x)\) 在 \([0,3]\) 上连续。由题设
\[
2f(0)=\int_0^2 f(x)\,dx,
\]
又
\[
2f(0)=\int_0^2 f(0)\,dx,
\]
所以
\[
\int_0^2 \bigl[f(x)-f(0)\bigr]dx=0,
\]
即
\[
\int_0^2 g(x)\,dx=0.
\]

若在 \((0,2)\) 内不存在 \(\eta\) 使 \(g(\eta)=0\)，则由连续性，\(g(x)\) 在 \((0,2)\) 内只能恒正或恒负。若 \(g(x)>0\)，则
\[
\int_0^2 g(x)\,dx>0,
\]
矛盾；若 \(g(x)<0\)，则
\[
\int_0^2 g(x)\,dx<0,
\]
也矛盾。因此存在 \(\eta\in(0,2)\)，使得
\[
g(\eta)=0,
\]
即
\[
f(\eta)=f(0).
\]

因为 \(f(x)\) 在 \([0,\eta]\) 上连续，在 \((0,\eta)\) 内可导，且
\[
f(0)=f(\eta),
\]
由罗尔定理，存在 \(\xi\in(0,\eta)\)，使得
\[
f'(\xi)=0.
\]
又 \(0<\eta<2<3\)，故 \(\xi\in(0,3)\)。因此
\[
\exists\,\eta\in(0,2),\ f(\eta)=f(0),\qquad \exists\,\xi\in(0,3),\ f'(\xi)=0.
\]
\end{solutioncontent}

\mistakeknowledge{
\begin{itemize}
  \item \textbf{核心公式/定理}：若连续函数 \(g\) 在区间内不取零，则不能变号；若 \(\int_a^b g(x)\,dx=0\)，则 \(g\) 不能在区间内恒正或恒负。罗尔定理：连续、可导且端点函数值相等，则存在 \(\xi\) 使 \(f'(\xi)=0\)。
  \item \textbf{易错点}：不能只用积分中值定理得到闭区间上的点就结束，题目要求 \(\eta\in(0,2)\)，需排除端点并证明开区间内存在零点。
  \item \textbf{关联知识}：由 \(2f(0)=\int_0^2 f(x)\,dx\) 联想到 \(2f(0)=\int_0^2 f(0)\,dx\)，从而构造 \(\int_0^2 [f(x)-f(0)]\,dx=0\)。
\end{itemize}}

\end{mistake}
